\chapter*{第 2 章 流体静力学}
\addcontentsline{toc}{chapter}{第 2 章 流体静力学}
\markboth{第 2 章 流体静力学}{}

\begin{tishi}
第 2 章是 818 计算题的\textbf{基础入门}：E1 slide33 明确"第 2 章流体静力学＝计算题基础入门、重点掌握方法"。
本章 17 题中，2.1--2.6 是测压计与测压管水头（静压强测量、压强表示法），2.7--2.9 是旋转容器相对平衡，
2.10--2.17 是平面与曲面上的总压力、压力中心（本章最可能出大题的部分）。
做题要求：计算题必须写出\textbf{已知、依据（公式号）、代入、结果（带单位）}四步；
曲面总压力题必须先画压力体并判断\textbf{实压力体还是虚压力体}，再定铅直分力的方向。
\end{tishi}

\begin{ti}
\sloppy
\emergencystretch=2em

\item 如习题 2.1 图所示的密封容器，容器盛的液体是汽油，$\gamma=7.35\ \mathrm{kN/m^3}$，当已知测压管高出液面 $h=1.5\ \mathrm{m}$，求用水柱高度表示的液面相对压强 $p_0$。
\tuyi{习题 2.1 图：左侧为密封容器，容器内盛汽油，液面处标出液面相对压强 $p_0$（待求）；容器右壁下部接一根竖直开口测压管（通大气），管内汽油液面高出容器内液面 $h=1.5\ \mathrm{m}$。}
\xstarfull{4}\yiju{E6 教材式 2.11、2.12（静压强测量与压强表示法，讲义定为四星考点）；E1 slide33"第 2 章＝计算题基础入门、重点掌握方法"；E4 2015 年试题填空题考"U 形水银压差计"读数换算，同属测压计类计算}\nandu{易}

\item 习题 2.2 图中压差计上部有空气，$h_1=0.6\ \mathrm{m}$，$h=0.45\ \mathrm{m}$，$h_2=1.8\ \mathrm{m}$，求 A、B 两点压强差，工作介质水的 $\gamma=9800\ \mathrm{N/m^3}$。
\tuyi{习题 2.2 图：倒 U 形压差计，左端接测点 A 的球形容器，右端接测点 B 的球形容器；压差计上部（含顶端阀门 C 的一段）充满\textbf{空气}，空气自重忽略不计；左腿液面 D 与右腿液面 E 之间铅直距离为 $h=0.45\ \mathrm{m}$；液面 D 到 A 接管轴线的铅直距离 $h_1=0.6\ \mathrm{m}$；A 接管轴线到 B 接管轴线的铅直距离 $h_2=1.8\ \mathrm{m}$。}
\xstarfull{4}\yiju{E4 2015 年试题填空题"U 形水银压差计"（E6 教材图 2.7 压差计公式、式 2.9、2.10）；E1 slide33"第 2 章＝重点掌握方法"；E2"第 1--3 章偏概念，掌握基本公式与解释即可"}\nandu{易}

\item 如习题 2.3 图所示为一复式水银测压计，用以测量水箱中水的表面相对压强。根据图中读数（单位为 m）计算水面相对压强值。
\tuyi{习题 2.3 图：复式水银测压计。水箱内水面（标 $p_0$）高程 3.0；水箱右侧经水管接第一支水银 U 形管，其两腿水银面高程分别为 1.4 与 2.5；高程 2.5 与 1.2 两个水银面之间由\textbf{充水的倒 U 形管}连通；第二支 U 形管左腿水银面高程 1.2、右腿（开口通大气）水银面高程 2.3。各读数单位均为 m，高程自水箱底量起。}
\xstarfull{4}\yiju{E6 教材图 2.7 复式测压计等压面法、式 2.9、2.10（静压强测量，讲义四星考点）；E4 2015 年填空考 U 形水银压差计；E1 slide33"第 2 章＝重点掌握方法"}\nandu{中}

\item 如习题 2.4 图所示，$h_1=0.5\ \mathrm{m}$，$h_2=1.8\ \mathrm{m}$，$h_3=1.2\ \mathrm{m}$，试根据水银压力计的读数，求水管 A 内的真空度及绝对压强。（设大气压的压力水头为 10 m）
\tuyi{习题 2.4 图：水管 A 内充满水；自 A 向右接一根先向上翻越、再向下插入水银的管子（管内充水），翻越后向下与水银接触，该封闭侧水银面高程 $h_2=1.8\ \mathrm{m}$；水银经 U 形底部上升到通大气的开口腿，开口侧水银面高程 $h_3=1.2\ \mathrm{m}$；水管 A 的轴线高程 $h_1=0.5\ \mathrm{m}$。三个高程均自图底地面基准线量起。}
\xstarfull{5}\yiju{E6 教材【例 2.1】与本题同型同数据（$h_1=0.5$、$h_2=1.8$、$h_3=1.2$，$p_a=98000\ \mathrm{Pa}$，答 $p_{A,\mathrm{abs}}=30772\ \mathrm{Pa}$、真空度 $67228\ \mathrm{Pa}$），是教材正文唯一用"真空度＋绝对压强"两步求解的例题；E5 题库选择题考"虹吸管最高处压强小于大气压"（真空度概念）；E1 slide33"第 2 章＝重点掌握方法"}\nandu{中}

\item 如习题 2.5 图所示，敞开容器内注有三种互不相混的液体，$\gamma_1=0.8\gamma_2$，$\gamma_2=0.8\gamma_3$，求侧壁处三根测压管内液面至容器底部的高度 $h_1$、$h_2$、$h_3$。
\tuyi{习题 2.5 图：敞开容器内自上而下注入三种互不相混的液体，每层厚 2 m，自上而下重度依次为 $\gamma_1$、$\gamma_2$、$\gamma_3$；容器右侧壁接三根测压管——第一根接在 $\gamma_1$ 与 $\gamma_2$ 交界面处（图中液面最高，标 $h_3$），第二根接在 $\gamma_2$ 与 $\gamma_3$ 交界面处（标 $h_2$），第三根接在容器底部 $\gamma_3$ 层内（图中液面最低，标 $h_1$）；三管内液面至容器底部的高度分别为 $h_1$、$h_2$、$h_3$。}
\xstarfull{4}\yiju{E6 教材式 2.9、2.10（测压管水头 $z+p/(\rho g)=C$，讲义"静力学基本方程与测压管水头"为五星考点）；E1 slide33"第 2 章＝计算题基础入门、重点掌握方法"；E2"第 1--3 章偏概念，掌握基本公式与解释即可"}\nandu{中}

\item 试标出习题 2.6 图示盛液容器内 A、B 和 C 三点的位置水头、压强水头和测压管水头。以习题 2.6 图示 O--O 为基准面。
\tuyi{习题 2.6 图：密闭盛液容器，液面处标出液面压强 $p_0$；容器内标出 A（下部偏右）、B（下部偏左）、C（上部居中）三点；容器右壁接一根向上开口通大气（标 $P_a$）的测压管，管内液面高于容器内液面；容器下方有一条水平基准线，标注 O--O。}
\xstarfull{5}\yiju{E6 教材式 2.9、2.10 及"静力学基本方程与测压管水头"（讲义定为五星考点）；E4 2022 年 818 单选第 3 题考"测压管水头线坡度小于 0 的沿程变化"、第 7 题考"虹吸管最高处压强与大气压的关系"；E5 题库选择题"虹吸管最高处压强小于大气压"}\nandu{易}

\item 一旋转圆柱容器高 $H=0.6\ \mathrm{m}$，直径 $D=0.45\ \mathrm{m}$，容器中盛水，如习题 2.7 图所示。求水正好与容器中心触底，顶部与容器同高时，容器的旋转角速度 $\omega$。
\tuyi{习题 2.7 图：圆筒形旋转容器，$z$ 为竖直旋转轴，角速度 $\omega$；容器高 $H=0.6\ \mathrm{m}$、直径 $D=0.45\ \mathrm{m}$（半径 $R=0.225\ \mathrm{m}$）；旋转后自由面为旋转抛物面，抛物面顶点恰好落在容器中心底部，而液面在容器壁处恰好升至容器顶面（与容器同高）。}
\xstarfull{3}\yiju{E6 教材式 2.16、2.17（自由面方程 $z_0=\omega^2r^2/(2g)$），讲义把"旋转容器"列为三星考点；E1 slide33"第 2 章＝重点掌握方法"；E2"第 1--3 章偏概念"}\nandu{易}

\item 有盖圆筒形容器半径为 $R$，高 $H$，绕垂直轴心线以恒角速度旋转，求筒中水压强表达式，分容器上盖中心开一小孔和顶盖边缘开口两种情况，当时当地大气压力为 $p_a$。
\xstarfull{3}\yiju{E6 教材式 2.14、2.15（$p=\rho\omega^2r^2/2-\rho gz+C$ 与边界条件），讲义"旋转容器"三星考点；E2"第 1--3 章偏概念，掌握基本公式与解释即可"；本题是推导题，历年真题中未单独成题，故压至三星}\nandu{中}

\item 如习题 2.9 图所示，已知容器的半径 $R=15\ \mathrm{cm}$，高 $H=50\ \mathrm{cm}$，充水深度 $h=30\ \mathrm{cm}$，若容器绕 $z$ 轴以等角速度 $\omega$ 旋转，试求：容器以多大极限转速旋转时，才不致使水从容器中溢出。
\tuyi{习题 2.9 图：圆筒形容器半径 $R=15\ \mathrm{cm}$、高 $H=50\ \mathrm{cm}$；绕竖直的 $z$ 轴以等角速度 $\omega$ 旋转；旋转前充水深度 $h=30\ \mathrm{cm}$（自容器底量起）；旋转后自由面为抛物面，轴线处液面下降、壁面处液面升高。}
\xstarfull{3}\yiju{E6 教材式 2.17 自由面方程与旋转前后体积不变，讲义"旋转容器"三星考点；E1 slide33"第 2 章＝重点掌握方法"；历年真题未单独成题，故定为三星}\nandu{中}

\item 在什么特殊情况下，水下平面的压力中心与平面形心重合？
\xstarfull{4}\yiju{E6 教材式 2.29--2.31（$y_D=y_C+J_C/(y_C A)$，讲义"平面总压力与压力中心"为五星考点）；E1 slide16"判断、选择题考概念辨析"；E2"第 1--3 章偏概念，掌握基本公式与解释即可"；本册前言已把 2.10 点名为必须写出完整答案的概念题}\nandu{易}

\item 如习题 2.11 图所示，一铅直矩形闸门，已知 $h_1=1\ \mathrm{m}$，$h_2=2\ \mathrm{m}$，宽 $b=1.5\ \mathrm{m}$，分别用解析公式和图解法求总压力及其作用点。
\tuyi{习题 2.11 图：铅直矩形闸门 AB（A 为上缘、B 为下缘，闸门铅直放置）；闸门顶缘 A 在水面以下 $h_1=1\ \mathrm{m}$ 处，闸门高 $h_2=2\ \mathrm{m}$、宽 $b=1.5\ \mathrm{m}$；水自左侧作用在闸门上（图中箭头 F 水平指向闸门）。}
\xstarfull{5}\yiju{E6 教材式 2.29--2.31 与【例 2.3】（铅直矩形闸门，解析法与图解法并列，图解法 $P=\Omega b$）；E4 2015 年试题计算题考"圆柱闸门静水总压力"；E5 题库计算题考"弧形闸门静水总压力 $F_x$、$F_z$ 及力矩"；讲义"平面总压力与压力中心"为五星考点}\nandu{中}

\item 蓄水池侧壁装有一直径为 $D$ 的圆形闸门，闸门平面与水面夹角为 $\theta$，闸门形心 C 处水深 $h_C$，闸门可绕通过形心 C 的水平轴旋转，如习题 2.12 图所示，证明作用于闸门水压力对轴的力矩与形心水深 $h_C$ 无关。
\tuyi{习题 2.12 图：蓄水池侧壁上装一倾斜圆形闸门，闸门平面与水面夹角 $\theta$；闸门直径 $D=125\ \mathrm{cm}$；图中标出上游水深 $h_0$、形心 C 处水深 $h_C$、压力中心 D 及总压力 P 的作用线；闸门可绕通过形心 C 的水平轴旋转。}
\xstarfull{5}\yiju{E6 教材式 2.29--2.31（$P=\rho gh_CA$、$y_D-y_C=J_C/(y_CA)$，圆形 $J_C=\pi D^4/64$）；E5 题库计算题考"弧形闸门静水总压力及力矩"；讲义"平面总压力与压力中心"五星考点；本题是教材正文压力中心公式的直接推论，属必考证明型}\nandu{中}

\item 如习题 2.13 图所示矩形闸门 AB 的宽 $b=3\ \mathrm{m}$，$\alpha=60^{\circ}$，$h_1=1\ \mathrm{m}$，$h_2=1.73\ \mathrm{m}$，$h_3=h_2/2$，闸门自重 $G=9800\ \mathrm{N}$，试求：将门吊起所需的力 $T$。
\tuyi{习题 2.13 图：矩形闸门 AB，宽 $b=3\ \mathrm{m}$（垂直于图面）；A 端铰接在右侧竖直壁上，闸门自 A 向左下倾斜、与底板成 $\alpha=60^{\circ}$，下缘 B 落在底板上；上游水面高出 A 点 $h_1=1\ \mathrm{m}$，A 到底板的铅直距离 $h_2=1.73\ \mathrm{m}$；闸门下游（右侧）水位高出底板 $h_3=h_2/2$；闸门自重 $G=9800\ \mathrm{N}$；图中竖直向上的箭头 $T$ 为吊索拉力，其竖直作用线在闸门下缘 B 附近通过。}
\xstarfull{5}\yiju{E6 教材式 2.29--2.31 及【例 2.3】方法（先求总压力与压力中心，再取铰链力矩平衡）；E4 2015 年计算题考"圆柱闸门静水总压力"、2026 年回忆版计算题仍以虹吸／闸门／管嘴类经典题为主（E4）；E1 slide33"第 2 章＝重点掌握方法"}\nandu{难}

\item 如习题 2.14 图所示，一弧形闸门 AB，宽 $b=4\ \mathrm{m}$，圆心角 $\alpha=45^{\circ}$，半径 $r=2\ \mathrm{m}$，闸门转轴恰与水面齐平，求作用于闸门的静水总压力。
\tuyi{习题 2.14 图：弧形闸门 AB，宽 $b=4\ \mathrm{m}$；闸门转轴（圆心 O）位于水面上，O 与 A 的水平距离为半径 $r=2\ \mathrm{m}$（A 在水面线上）；B 点在底部，$\angle AOB=\alpha=45^{\circ}$；水深 $h=r\sin\alpha=1.414\ \mathrm{m}$；水在闸门左侧。}
\xstarfull{5}\yiju{E5 题库计算题"弧形闸门静水总压力 $F_x$、$F_z$ 及力矩"；E4 2015 年计算题考"圆柱闸门静水总压力"；E6 教材式 2.32--2.35 与【例 2.4】（坝顶圆柱形闸门，弧面压力体与合力过圆心）；讲义"曲面总压力""压力体虚实"均为五星考点}\nandu{中}

\item 如习题 2.15 图所示，圆柱闸门长 $L=4\ \mathrm{m}$，直径 $D=1\ \mathrm{m}$，上下游水深分别为 $H_1=1\ \mathrm{m}$，$H_2=0.5\ \mathrm{m}$，试求长柱体上所受的静水总压力。
\tuyi{习题 2.15 图：直径 $D=1\ \mathrm{m}$ 的圆柱闸门横卧在底板上，长 $L=4\ \mathrm{m}$（垂直于图面）；左侧上游水深 $H_1=1\ \mathrm{m}$（自底板量起，恰与圆柱顶部齐平），右侧下游水深 $H_2=0.5\ \mathrm{m}$（自底板量起，恰与圆柱轴线齐平）。}
\xstarfull{5}\yiju{E4 2015 年试题计算题"圆柱闸门静水总压力"（直接对应）；E6 教材式 2.32--2.35、【例 2.4】与【例 2.5】（曲面总压力分两侧、压力体分解为半圆与四分之一圆）；讲义"曲面总压力""压力体虚实"五星考点}\nandu{中}

\item 一水坝受水面为一抛物线，顶点在 O 点，水深 $H=50\ \mathrm{m}$ 处抛物线到抛物线对称轴距离为 12.5 m，如习题 2.16 图所示，求水作用于单位宽度坝体的合力大小和方向。
\tuyi{习题 2.16 图：坝的上游受水面为抛物线，抛物线顶点 O 位于坝底、对称轴为过 O 的竖直点划线；水面高程为 $H=50\ \mathrm{m}$，水面上 A 点到对称轴的水平距离为 12.5 m（即抛物线过点 $(12.5,50)$）；图下方为坝体与地基；要求单位宽度（垂直于图面 1 m）坝体上水的合力。}
\xstarfull{5}\yiju{E6 教材式 2.32--2.35（曲面总压力：$P_x=\rho gh_CA_x$、$P_z=\rho gV$）与【例 2.4】【例 2.5】的曲面积分／压力体方法；E5 题库计算题考"弧形闸门静水总压力"；讲义"曲面总压力""压力体虚实"五星考点；需先由"顶点在 O、过 $(12.5,50)$"定出抛物线方程再积分，故难度较大}\nandu{难}

\item 如习题 2.17 图所示，扇形闸门，中心角 $\alpha=45^{\circ}$，宽度 $B=1\ \mathrm{m}$（垂直于图面），可以绕铰链 C 旋转，用以蓄水或泄水。水深 $H=3\ \mathrm{m}$，确定水作用于此闸门上的总压力 $p$ 的大小和方向。
\tuyi{习题 2.17 图：扇形闸门可绕铰链 C 旋转，C 位于底板高程上；弧面各点到 C 的距离均为半径 $r$；过 C 的水平线与过 C 向左上方 $45^{\circ}$ 的半径分别交弧面于 $a$、$b$；$b$ 点位于水面上（水深 $H=3\ \mathrm{m}$），$a$ 点位于底板上；闸门宽度 $B=1\ \mathrm{m}$（垂直于图面）；水在闸门左侧，只有 $b$ 到 $a$ 的一段弧面过水。}
\xstarfull{5}\yiju{E5 题库计算题"弧形闸门静水总压力 $F_x$、$F_z$ 及力矩"；E6 教材式 2.32--2.35 与【例 2.4】（曲面各点压力沿法线过圆心，故合力过圆心）；讲义"曲面总压力""压力体虚实"五星考点}\nandu{难}

\end{ti}
